
\documentclass[12pt]{article}
\usepackage[a4paper, total={6.5in, 9.5in}]{geometry}
\usepackage{graphicx}			
\usepackage{amssymb}
\usepackage{amsmath}			
\usepackage{fancyhdr}
\usepackage{enumerate}       
\usepackage{dirtytalk} 
\usepackage[english]{babel}
\usepackage{amsthm}
\usepackage{xcolor}
\usepackage[shortlabels]{enumitem}
\usepackage{titlesec}
\titleformat{\section}{\normalfont\Large\bfseries}{\thesection}{0pt}{\Large}
\setlength{\parindent}{0pt}
\setlength{\parskip}{0.8em}

\newtheorem{theorem}{Theorem}[section]
\newtheorem{corollary}{Corollary}[theorem]
\newtheorem{lemma}[theorem]{Lemma}
\newtheorem{definition}[theorem]{Definition}


\begin{document}
% \renewcommand{\thesection}{Exercise \arabic{section}}


\setcounter{section}{0} % Sets the section counter to 4, so the next section is 5
\renewcommand{\thesection}{}


\parindent = 0pt


      {\bf Analysis 1 HW 4}
%==================

\section{2.6}
{\bf Definitions}\\
$E\subset X$. \\
$E'$ is the set of all limit points of $E$.\\
$\bar{E}=E\cup E'$.\\

{\bf Prove that $E'$ is closed.}
\begin{proof}
To show $E'$ is closed, we will show that $E'$ contains all of its 
limit points. Let $p$ be a limit point $E'$. We will show that $p$ 
is a limit point of $E$, which in turn implies $p\in E'$. Let $\epsilon>0$ 
be arbitrary. $B_{\frac{\epsilon}{2}}(p)\cap E'$ contains at least one 
other point $q$ other than $p$. Also, since $q$ is a limit point of $E$, 
$B_{\frac{\epsilon}{2}}(q)\cap E$ contains at least one point $r$ other 
than $q$. 
\begin{align*}
      d(p, r) & \le d(p, q)+d(q, r) \\
       & < \frac{\epsilon}{2} + \frac{\epsilon}{2}=\epsilon.
\end{align*}
      
So, for every $\epsilon>0$, we have $r\in B_{\epsilon}(p)\cap E$, $r\ne p$. 
Therefore, $p$ is a limit point of $E$.
\end{proof}

{\bf Prove that $E$ and $\bar{E}$ have the same limit points.}
\begin{proof}
It is clear that if $p$ is a limit point of $E$, it must also be a limit 
point of $\bar{E}$. Going the other way, let $p$ be a limit point of 
$\bar{E}$. Let $\epsilon>0$ be arbitrary. Then, 
$B_{\frac{\epsilon}{2}}(p)\cap \bar{E}$ contains at least one $q$ other 
than $p$. If $q\in E$, we are done. If $q\not\in E$, it must be the case 
that $q\in E'$, i.e, $q$ is a limit point of $E$. Thus, 
$B_{\frac{\epsilon}{2}}(q)\cap E$ must contain a point $r$ other than 
$q$. Now, $d(p, r)\le d(p, q)+d(q, r)< \frac{\epsilon}{2} +
\frac{\epsilon}{2}=\epsilon$. So, $E$ and $\bar{E}$ have the same limit 
points.
\end{proof}
{\bf Do $E$ and $E'$ always have the same limit points?}

No. Consider $E=\left\{  \frac{1}{n}\ | \ n\in \mathbb{N}  \right\}$. 
$E'=\{ 0 \}$. The set of limit points of $E'$ is the null set. 


\section{2.29}

{\bf Prove that every open set in $\mathbb{R}^{1}$ is a union of an at 
most countable collection of disjoint segments.}

\begin{proof}
      Let $U$ be an open set in $\mathbb{R}$. Let $x, y, z\in U$. 
      We define $x\sim y$ if ($x\leq y$ and $[x, y]\subset U$) or ($y<x$ 
      and $[y, x]\subset U$). This relation is
      \begin{enumerate}
            \item reflexive, since $x\sim x$ ($[x, x]\subset U$),
            \item symmetric, since $x\sim y \implies y \sim x$, and
            \item transitive, since $x\sim y$ and $y\sim z$ $\implies$ $x\sim z$.
      \end{enumerate}
      Thus, this relation induces a partition $P$ on $U$. We know that all 
      elements of a partition are disjoint. To show that every $Q\in P$ is 
      a segment, consider $p\in Q$. Since $U$ is open, $r>0$ exists such 
      that $B_{r}(p)\subset U$. For every $b\in B_{r}(p)$, if $b\leq r$ we 
      have $[b, r]\subset B_{r}(p)\subset U$, and if $b>r$, we have 
      $[r, b]\subset B_{r}(p)\subset U$.So, $b\sim p$, i.e, $b$ and $p$ 
      must belong to the same equivalence class, $Q$. It follows that 
      $B_{r}(p)\subset Q$. Therefore, each $Q\in P$ much be a segment 
      in $\mathbb{R}$. 

      Now, for every $Q_{i}\in P$, define 
      $q_{i}\equiv({\inf Q_{i}+\sup Q_{i}})/2$. Since $Q_{i}$ are disjoint, 
      each $Q_{i}$ has a unique $q_{i}$ associated with it. Thus, 
      we have $P\sim \mathbb{Q}'\subset \mathbb{Q}$. Since $\mathbb{Q}$ 
      is countable, $P$ is at most countable. 
\end{proof}



\section{4.2}

{\bf Definitions}\\
$f:X\to Y$ is continuous, where $X$ and $Y$ are metric spaces.\\

{\bf Prove that $f(\overline{E})\subset\overline{f(E)}$ for every set $E\subset X$. }


\begin{proof}
      Let $p \in f(\overline{E})$. Then, there exists a 
      $q$ in $\overline{E}$ such that $f(q)=p$. If $q\in E$, we are 
      done. If $q\not\in E$, $q$ must be a limit point of $E$. So, there 
      exists a sequence $(q_{n})$ in $E$ which converges to $q$. Since 
      $f$ is continuous, $(f(q_{n}))$ must converge to $f(q)$. Thus, 
      $f(q)$ is a limit point of $f(E)$, and so $f(q)=p\in \overline{f(E)}$. 
\end{proof}

{\bf Show, by an example, that $f(\overline E)$ can be a proper subset of $\overline{f(E)}$.}

Consider $f:\mathbb{N}\to \mathbb{R}$, $f(n)=\frac{1}{n}$. Since 
$\mathbb{N}$ consists of only isolated points, $f$ is continuous. Also, 
$\overline{\mathbb{N}}=\mathbb{N}$, since $\mathbb{N}$ is a closed set. 
Thus, $f(\mathbb{N})= f(\overline{\mathbb{N}})=\left\{  \frac{1}{n} \ 
| \ n\in \mathbb{N} \right\}$. However, $0$ is a limit point of 
$f(\mathbb{N})$. Thus, $0\in \overline{f(\mathbb{N})}$, but 
$0\not\in f(\overline{\mathbb{N}})$. 

\section{4.4}
{\bf Definitions}\\
$f, g:X\to Y$ are continuous, where $X$ and $Y$ are metric spaces.\\
$E\subset X$ is dense in $X$.\\

{\bf Prove that $f(E)$ is dense in $f(X)$.}

\begin{proof}
      Let $p\in f(X)$. There exists $q\in X$ such that $f(q)=p$. 
      Since $E$ is dense in $X$, either $q\in E$ or $q$ is a limit point 
      of $E$. If $q\in E$, $f(q)\in f(E)$. If $q$ is a limit point of $E$, 
      there exists a sequence $(q_{n})$ in $E$ which converges to $q$. 
      Since $f$ is continuous, $(f(q_{n}))$ must converge to $f(q)$. Thus, 
      $f(q)$ is a limit point of $f(E)$. Therefore, either $p\in f(E)$, or 
      $p$ is a limit point of $f(E)$, i.e, $f(E)$ is dense in $f(X)$. 
\end{proof}

{\bf 
If $g(p)=f(p)$ for all $p\in E$, prove that $g(p)=f(p)$ for all $p\in X$. }

\begin{proof}
      Again, $E$ being dense in $X$ tells us that either $p\in E$ or $p$ is 
      a limit point of $E$. If $p\in E$, the result is true by hypothesis. 
      In the latter case, there exists a sequence $(p_{n})$ in $E$ that 
      converges to $p$. Since $f$ and $g$ are continuous, 
      $(f(p_{n}))\to f(p)$ and $(g(p_{n}))\to g(p)$. Also, 
      $f(p_{n})=g(p_{n})$ for all $n\in \mathbb{N}$, by hypothesis. 
      Since the limit of a sequence is unique, it follows that $f(p)=g(p)$.  
\end{proof}


\section{4.7}
{\bf Definitions}\\
$f, g:\mathbb{R}^{2}\to \mathbb{R}$.\\\\
$
f(x, y)=\begin{cases}
0 &  (x, y)=(0, 0)\\
\frac{xy^{2}}{x^{2}+y^{4}} & (x, y)\ne(0, 0)
\end{cases}
$

$
g(x, y)=\begin{cases}
0 &  (x, y)=(0, 0)\\
\frac{xy^{2}}{x^{2}+y^{6}} & (x, y)\ne(0, 0)
\end{cases}
$\\

{\bf 
Prove that $f$ is bounded on $\mathbb{R}^{2}$.}

\begin{proof}
      We will show that
$$
\left|\frac{xy^{2}}{x^{2}+y^{4}}\right|\leq \frac{1}{2}
$$
for all $(x, y)\in \mathbb{R}^{2}$. On applying the AM-GM inequality on $x^{2}$ and $y^{4}$, we get 
\begin{align*}
\frac{x^{2} +y^{4}}{2}\geq \sqrt{ x^{2}y^{4} }=|x|y^{2} \\
\implies \frac{|x|y^{2}}{x^{2}+y^{4}}=\left|\frac{xy^{2}}{x^{2}+y^{4}}\right|\leq \frac{1}{2}.
\end{align*}
\end{proof}

{\bf 
Prove that $g$ is unbounded on every neighborhood of $(0, 0)$.}

\begin{proof}
      Let $\epsilon>0$. We will show that there exists $(x, y)\in B_{\epsilon}((0, 0), \mathbb{R}^{2})$ such that
\begin{align}
|g(x, y)|=\left|\frac{xy^{2}}{x^{2}+y^{6}}\right|>M
\end{align}
for all $M>0$. Assume $x, y>0$. 
\begin{align*}
\frac{xy^{2}}{x^{2}+y^{6}}>M \\
\implies \frac{{x^{2}+y^{6}}}{xy^{2}}=\frac{x}{y^{2}}+\frac{y^{4}}{x}< \frac{1}{M}
\end{align*}
Define $\alpha\equiv x/y^{2}$.
\begin{align*}
\alpha+\frac{y^{2}}{\alpha}< \frac{1}{M}.
\end{align*}
Let $\alpha=\frac{1}{2M}$. Let us attempt to solve for $y$.
\begin{align*}
\frac{1}{2M} +2My^{2}&< \frac{1}{M} \\
y&<\frac{1}{2M}
\end{align*}
Setting $y=\frac{1}{4M}$ will suffice. Now, we can solve for $x$:
\begin{align*}
x=\frac{y^{2}}{2M}=\frac{1}{32M^{3}}
\end{align*}
Thus, $\left( \frac{1}{32M^{3}}, \frac{1}{4M} \right)$ will satisfy (1). This can be easily verified:
\begin{align*}
g\left( \frac{1}{32M^{3}}, \frac{1}{4M} \right)={{\frac{{{{{{{\left({{{\frac{1}{{{4M}}}}}}\right)}}^{2}}}{\left({{{\frac{1}{{{{{32{M}^{3}}}}}}}}}\right)}}}}{{{{{{{\left({{{\frac{1}{{{4M}}}}}}\right)}}^{6}}}}+{{{{{\left({{{\frac{1}{{{{{32{M}^{3}}}}}}}}}\right)}}^{2}}}}}}}} = \frac{8M}{5}>M.
\end{align*}
Now, if $|x|< \epsilon/\sqrt{ 2 }$ and $|y|<\epsilon/\sqrt{ 2 }$, it must be that $(x, y)\in B_{\epsilon}(0, 0)$. There exist $k_{1}\geq1$ and $k_{2}\geq 1$ such that $\frac{1}{32(k_{1}M)^{3}}< \frac{\epsilon}{\sqrt{ 2 }}$ and $\frac{1}{4k_{2}M}< \frac{\epsilon}{\sqrt{ 2 }}$. Let $k$ be the greater of the two. Since being greater than $kM$ ensures being greater than $M$, we have found the $(x, y)$ we were in the search of:
\begin{align*}
(x, y)=\left( \frac{1}{32(kM)^{3}}, \frac{1}{4kM} \right).
\end{align*}
\end{proof}

{\bf 
Prove that $f$ is not continuous at $(0, 0)$.}

\begin{proof}
      Since $(0, 0)$ is a limit point of $\mathbb{R}^{2}$, we must have $\lim_{ (x, y) \to (0, 0) }f(x, y)=f(0, 0)=0$ for $f$ to be continuous at $(0, 0)$. If we plug $x=y^{2}$ into the definition of $f$, we get $\frac{1}{2}$. Thus, for $\epsilon< \frac{1}{2}$, it is not possible to find a $\delta>0$ such that $|(x, y)|<\delta \implies |f(x, y)|<\epsilon$, since we will always be able to find $(a, b)\in B_{\delta}(0, 0)$ such that $a=b^{2}$. 
\end{proof}








% \section{hello}
%  {\bf Definitions}\\
% $(x_{n})$ is a sequence of real numbers satisfying $3x_{n+1} = {x_{n}^{3}-2}$ for all $n\in \mathbb{N}$.

%       {\bf Solution}\\
% If $(x_{n})\to l$ for some $l$, then
% \begin{align*}
%                & \lim_{ n \to \infty } (3x_{n+1}) = \lim_{ n \to \infty } (x_{n}^{3}-2) \\
%       \implies & 3l=l^{3}-2
% \end{align*}
% courtesy the algebraic limit theorem. This yields $l\in \{ 2, -1 \}$.

% Now, consider the successive differences:
% \begin{align}
%       x_{n+1}-x_{n} & = \frac{{x_{n}^{3}-3x_{n}-2}}{3}
% \end{align}

% \begin{enumerate}[(i)]
%       \item If $x_{i} > 2$, $x_{i}^{3}-3x_{i}-2 = x_{i}(x_{i}^{2}-3)-2>2(4-3)-2=0$. So, $x_{i+1}-x_{i}>0$. Therefore, if $x_1>2$, the sequence is increasing and cannot converge to 2 or -1, so it must diverge.
%       \item If $x_{i} = 2$, $x_{i+1} -x_{i}=0$. Thus, if $x_1 =2$ the sequence is constant at 2, and it converges to 2.
%       \item If $-1<x_{i}<2$, we make the following claims:
%             \begin{enumerate}
%                   \item {\bf Claim: }If $-1<x_{i}<2$, then $x_{i+1}-x_{i}\leq0$.\\
%                         It suffices to show that $y^{3}-3y-2\leq0$ for all $y<2$.
%                         Consider the expression $t(6-t)-9$, $t>0$. If $t>6$, it is clear that $t(6-t)-9 <0$. If $0<t\le 6$ , $t$ and $6-t$ are non negative. From the AM-GM inequality, we have $9\le t(6-t) \implies$  $t(6-t)-9\le 0$. Thus, $t(t(6-t)-9)\le 0$ is true for all $t>0$. Now, let $y=2-t$. Plugging $t=2-y$ into the above inequality, we get ${{{y}^{3}}-{3y}-2}\le 0$ for all $y<2$.
%                   \item {\bf Claim: }If $-1<x_{i}<2$, then $-1<x_{i+1}<2$.\\
%                         \begin{equation*}
%                               x_{i}<2
%                               \implies x_{i}^{3}<8
%                               \implies  x_{i}^{3}-2<6
%                               \implies  \frac{{x_{i}^{3}-2}}{3}=x_{i+1}<2
%                         \end{equation*}

%                         \begin{equation*}
%                               -1<x_{i}
%                               \implies  -1<x_{i}^{3}
%                               \implies  -3<x_{i}^{3}-2
%                               \implies  -1< \frac{{x_{i}^{3}-2}}{3} = x_{i+1}
%                         \end{equation*}
%             \end{enumerate}
%             Thus, if $x_{1}$ lies between $-1$ and $2$, then $(x_{n})$ is bounded and monotone decreasing. By the monotone convergence theorem, the sequence must converge, and since it cannot converge to 2, it must converge to -1.
%       \item If $x_{i} = -1$, $x_{i+1} -x_{i}=0$. Thus, if $x_1 = -1$ the resulting sequence is constant at -1, and it converges to -1.
%       \item For $x_i < -1$, we have $x_i^3 - 3x_i - 2 \le 0$ from (iii)(a) (while claimed for $-1<x_i<2$, we proved it for all $x_i<2$). Thus, for $x_1<-1$ the sequence is decreasing. Since the sequence cannot converge to either -1 or 2, it must diverge to negative infinity.
% \end{enumerate}

% \section{}
%  {\bf Definitions}\\
% $\{ (a_{i}, b_{i}) \}$ is a sequence in $\mathbb{R}^{2}$ with $a_{i}\le b_{i}$ such that for each $i$, the closed interval $[a_{i}, b_{i}]$ contains the closed interval $[a_{i+1}, b_{i+1}]$.\\
% $\sup a_{n} = \sup \{ a_{k}\ |\ k\ge n\}$.

%       {\bf Solution}
% \begin{enumerate}[(a)]
%       \item {\bf To prove: }$\bigcap_{i}[a_{i}, b_{i}]\ne \emptyset$.
%             \begin{proof}
%                   Consider the sequences $(a_{n})$ and $(b_{n})$ independently.
%                   Since $[a_{i}, b_{i}]$ always contains $[a_{i+1}, b_{i+1}]$, $(a_{n})$
%                   must be monotone increasing, and $(b_{n})$ must be monotone decreasing.
%                   Also, every term of $(b_{n})$ is an upper bound for $(a_{n})$, and
%                   every term of $(a_{n})$ is a lower bound for $(b_{n})$. So,
%                   $\alpha = \sup a_{n}$ and $\beta = \inf b_{n}$ must exist.\\

%                   {\bf Claim: }$\alpha\leq \beta$.\\
%                   FTSOC, assume $\beta<\alpha$. From the definition of supremum,
%                   there must exist an $a_{i_{1}}$ such that $\beta<a_{i_{1}}<\alpha$.
%                   From the definition of infimum, there must exist a $b_{i_{2}}$ such
%                   that $\beta <b_{i_{2}}<a_{i_{1}}<\alpha$. This contradicts the fact that
%                   every term of $(b_{n})$ is an upper bound for $(a_{n})$.\hfill $\Rightarrow\Leftarrow$\\

%                   Thus, the interval $[\alpha, \beta]$ is always non empty.\\

%                   {\bf Claim: }$\bigcap_{i}[a_{i}, b_{i}]= [\alpha, \beta]$\\
%                   Every element in $[\alpha, \beta]$ is an upper bound for $(a_{i})$,
%                   and a lower bound for $(b_{i})$. Thus, $[\alpha, \beta]\subset[a_{i}, b_{i}]$
%                   for all $i$, i.e, $[\alpha, \beta]\subset \bigcap_{i}[a_{i}, b_{i}]$.\\
%                   If $k\in \bigcap_{i}[a_{i}, b_{i}]$, $k$ is an upper bound of $(a_{n})$
%                   and $k$ is a lower bound for $(b_{n})$. So, $k\ge \alpha$ and $k\le \beta$.
%                   Thus, $k \in [\alpha, \beta]$.\\

%                   Thus, $\bigcap_{i}[a_{i}, b_{i}]= [\alpha, \beta]\ne \emptyset$.
%             \end{proof}

%       \item {\bf To prove: }$|\bigcap_{i}[a_{i}, b_{i}]| =1\iff\lim_{ n \to \infty }(a_{n}-b_{n}) = 0$.
%             \begin{proof}
%                   We know that if a sequence is bounded and increasing,
%                   it must converge to the least upper bound of its range.
%                   Thus, $(a_{n})\to \alpha$. Similarly, if a sequence is bounded
%                   and decreasing, it must converge to the greatest lower bound of its
%                   range. Thus, $(b_{n})\to \beta$. Thus, we can say
%                   \begin{align*}
%                              & |\bigcap_{i}[a_{i}, b_{i}]| =1         \\
%                         \iff & \alpha=\beta                           \\
%                         \iff & \lim_{ n \to \infty }(a_{n}-b_{n}) = 0
%                   \end{align*}
%             \end{proof}
% \end{enumerate}
% \section{}
%  {\bf Definitions}\\
% $(a_{n})$ and $(b_{n})$ are sequences in $\mathbb{R}$.\\
% $\sup a_{n} = \sup \{ a_{k}\ |\ k\ge n\}$.


%       {\bf Solution}

% \begin{lemma}$\sup(a_{n}+b_{n})\le \sup a_{n}+ \sup b_{n}$ for all $n\in \mathbb{N}$\end{lemma}
% \begin{proof}
%       If $\sup (a_n)$ or $\sup (b_n)$ is $\infty$, the inequality is trivially true.
%       Thus, assume $\sup(a_n)$, $\sup(b_n) \in \mathbb{R}$.\\
%       FTSOC, assume $\sup a_{n} +\sup b_{n}<\sup(a_{n}+b_{n})$.
%       There must exist $a_{i}+ b_{i},\  i\ge n$ such that
%       $\sup a_{n} +\sup b_{n}<a_{i}+b_{i}<\sup(a_{n}+b_{n})$. We know
%       that $\sup a_{n}\ge a_{k}$ for all $k\ge n$. So, we have
%       \begin{align*}
%                      & (\sup a_{n}-a_{i}) +\sup b_{n}<b_{i} \\
%             \implies & \sup b_{n}<b_{i},
%       \end{align*}
%       which is impossible.\hfill $\Rightarrow\Leftarrow$\\
% \end{proof}

% \begin{lemma}If $(a_{n})\to a\in \mathbb{R}\cup \{ \infty, -\infty \}$, $(b_{n})\to b\in \mathbb{R}\cup \{ \infty, -\infty \}$, and $a_{n}\le b_{n}$ for all $n\in \mathbb{N}$, then $a\le b$. \end{lemma}

% \begin{proof}
%       If $(b_{n})\to \infty$, the inequality is trivially true.
%       If $(a_{n}) \to \infty$, it must be the case that $(b_{n})\to \infty$. Thus, assume
%       $a, b\in \mathbb{R}$. \\
%       FTSOC, assume $b<a$. Let $\epsilon = \frac{{a-b}}{2}$. Then, there exists an
%       integer $N_{1}$ such that for all $n\ge N_{1}$, $|a_{n}-a|<\epsilon$. Similarly,
%       there exists an integer $N_{2}$ such that for all $n\ge N_{2}$, $|b_{n}-b|<\epsilon$.
%       Let $N$ be the greater of $N_{1}$ and $N_{2}$. Thus, for all
%       $n\ge N$, $b_{n} \in \left( \frac{{3b-a}}{2}, \frac{{a+b}}{2} \right)$ and
%       $a_{n}\in \left( \frac{{a+b}}{2}, \frac{{3a-b}}{2} \right)$, which implies
%       $a_{n}> b_{n}$ for all $n\ge N$. \hfill $\Rightarrow\Leftarrow$\\
% \end{proof}

% {\bf To prove: }$\limsup_{ n \to \infty }(a_{n}+ b_{n})\le \limsup_{ n \to \infty}a_{n}+\limsup_{ n \to \infty } b_{n}$,
% provided the sum on the right is not of the form $\infty-\infty$.

% \begin{proof}
%       Consider the sequences $(\sup a_{n}+\sup b_{n})$, and $(\sup(a_{n}+b_{n}))$. From Lemma 3.1,
%       we have $\sup(a_{n}+b_{n})\le \sup a_{n}+ \sup b_{n}$ for all $n\in \mathbb{N}$.
%       Also, it is clear the the suprema of tails of a sequence form a monotone decreasing sequence.
%       Since all sequences are bounded in the extended reals, it follows these sequences must
%       converge to a value in the extended reals. Thus, we can use Lemma 3.2:
%       \begin{align*}
%                      & \sup(a_{n}+b_{n})\le \sup a_{n}+ \sup b_{n}                                                                                                                          \\
%             \implies & \lim_{ n \to \infty } (\sup(a_{n}+b_{n}))\le \lim_{ n \to \infty } (\sup a_{n}+ \sup b_{n}) = \lim_{ n \to \infty } (\sup a_{n})+ \lim_{ n \to \infty } (\sup b_{n}) \\
%             \implies & \limsup_{ n \to \infty }(a_{n}+ b_{n})\le \limsup_{ n \to \infty}a_{n}+\limsup_{ n \to \infty } b_{n}
%       \end{align*}
% \end{proof}

% \section{}
%  {\bf Definitions}\\
% $(p_{n})$ and $(q_{n})$ are Cauchy sequences in a metric space $(X, d)$.

%       {\bf Solution}

%       {\bf To prove: }The sequence $(d(p_{n}, q_{n}))$ converges.
% \begin{proof}

%       For any $m, n$, from the triangle inequality, we have
%       \begin{align*}
%             d(p_{n}, q_{n})\le d(p_{n}, p_{m}) + d(p_{m}, q_{n})
%       \end{align*}
%       and
%       \begin{align*}
%             d(p_{m}, q_{n})\le d(p_{m}, q_{m})+d(q_{m}, q_{n}).
%       \end{align*}

%       Combining the two, we get
%       \begin{align*}
%             d(p_{n}, q_{n})                            & \le d(p_{n}, p_{m}) + d(p_{m}, q_{m})+d(q_{m}, q_{n}) \\
%             \implies |d(p_{n}, q_{n})-d(p_{m}, q_{m})| & \le d(p_{n}, p_{m})+d(q_{m}, q_{n}).
%       \end{align*}
%       Let $\epsilon>0$ be arbitrary. Since $(p_{n})$ and $(q_{n})$ are Cauchy sequences, there must exist integers $N_{1}$ and $N_{2}$ such that for all $m, n\geq N_{1}$, $d(p_{n}, p_{m})< \frac{\epsilon}{2}$ and for all $m, n\geq N_{2}$, $d(q_{n}, q_{m})< \frac{\epsilon}{2}$ respectively. Let $N$ be the larger of $N_{1}$ and $N_{2}$. Thus,
%       \begin{align*}
%             |d(p_{n}, q_{n})-d(p_{m}, q_{m})| < \frac{\epsilon}{2} + \frac{\epsilon}{2} = \epsilon
%       \end{align*}
%       for all $n\ge N$. This implies that the sequence $(d(p_{n}, q_{n}))$ is Cauchy. In $\mathbb{R}$, a sequence is Cauchy if and only if it is convergent (to a limit in $\mathbb{R}$).
% \end{proof}

% \section{}
%  {\bf Definitions}\\
% $(p_{n})$ is a sequence in $\mathbb{R}$. \\
% For $a<b$, let $\zeta_{a, b}$ be the $[a, b]$-crossing number of $(p_{n})$.

%       {\bf Solution}

%       {\bf To prove: } $\zeta_{a, b}\in \mathbb{N}_{0} \ \ \forall \ a<b \iff (p_{n})\to p\in \mathbb{R}\cup \{ \infty, -\infty \}$.


% \begin{proof}
%       We will call a sequence with the property on the left (having finite $\zeta_{a,b}$ for all $a<b$) a \say{cross-bounded} sequence.

%             {\bf Proof of convergent $\implies$ cross-bounded}

%       We will prove the contrapositive: not cross-bounded $\implies$ not convergent.
%       If a sequence is not cross-bounded, $\zeta_{a, b}=\infty$ for some $a<b$. This implies
%       that for all $N$, there exists $d_{n}>N$ such that $p_{d_{n}}<a$,  and there exists $u_{n}>N$
%       such that $p_{u_{n}}>b$. This allows us to make the following deductions:
%       \begin{enumerate}[(i)]
%             \item For $\epsilon=b-a$, there does not exist an $N$ such that for all $m, n\ge N$, $|p_{m}-p_{n}|<\epsilon$.
%                   This tells us that $(p_{n})$ is not Cauchy, which implies $(p_{n})$ is not convergent in $\mathbb{R}$.
%             \item For all $M>a$, there does not exist an $N$ such that for all $n\geq N$, $p_{n}\geq M$. Thus, $(p_{n})$ cannot converge to $\infty$.
%             \item For all $m < b$, there does not exist an $N$ such that for all $n\geq N$, $p_{n}\leq m$. Thus, $(p_{n})$ cannot converge to $-\infty$.
%       \end{enumerate}

%       {\bf Proof of cross-bounded $\implies$ convergent}

%       We can rephrase the notion of being cross-bounded like so: for every $a<b$, there
%       either exists an $N$ such that for all $n\geq N$,  $p_{n}\ge a$, or there exists
%       an $N$ such that for all $n\ge N$, $p_{n}\leq b$. In other words, eventually $a$
%       becomes a lower bound for tails of $(p_n)$, or $b$ becomes an upper bound
%       for tails of $(p_n)$. We will analyze bounded (in $\mathbb{R}$) and unbounded
%       cases for $(p_{n})$ separately.
%       \begin{enumerate}[(i)]
%             \item Say $(p_{n})$ is bounded in $\mathbb{R}$, i.e, $m<p_{n}<M$ for all
%                   $n\in \mathbb{N}$ for some $m, M\in \mathbb{R}$. Consider the interval
%                   $$
%                         (\alpha, \beta):=\left(\frac{{3m+M}}{4}, \frac{{m+3M}}{4}\right).
%                   $$

%                   Given our paraphrasing of what it means to be cross-bounded, there
%                   either exists an $N$ such that for all $n\ge N$, $p_{n}\ge \alpha$, or
%                   there exists an $N$ such that for all $n\ge N$, $p_{n} \leq \beta$. In both
%                   cases, we have achieved new bounds for $(p_{n})$ for $n\ge N$, namely
%                   $(\alpha, M)$ and $(m, \beta)$ respectively. Notice that
%                   $M-a = \beta-m = \frac{3}{4}(M-m)$. Effectively, we have bounded a
%                   tail of the sequence in an interval that is three-fourths the size of the
%                   interval we started with. If we repeat this process $k$ times (always taking
%                   the maximum of all the $N$'s we've chosen), we will bound a tail of
%                   $(p_{n})$ in an interval of the size $\left( \frac{3}{4} \right)^{k}(M-m)$.
%                   Taking the limit as $k\to \infty$,
%                   $$
%                         \lim_{ k \to \infty } (M-m)\left( \frac{3}{4} \right)^{k} = (M-m)\lim_{ k \to \infty }\left( \frac{3}{4} \right)^{k} = 0
%                   $$
%                   Thus, for arbitrary $\epsilon>0$, there will exist $k$ such
%                   that $(M-m) \left( \frac{3}{4} \right)^{k}<\epsilon$, i.e, there will exist
%                   a tail of $(p_{n})$ where the difference between any two terms is less
%                   than $\epsilon$. This implies that $(p_{n})$ is Cauchy, and hence, convergent.
%             \item If $(p_{n})$ is not bounded in $\mathbb{R}$, we will show that it converges
%                   to $+\infty$ or $-\infty$. Consider any interval $(\alpha, \beta)$. As stated
%                   previously, one of these should be true:
%                   \begin{enumerate}[]
%                         \item There exists an $N$ such that for all $n\ge N$, $p_{n} \geq \alpha$.

%                               Since $(p_{n})$ is unbounded in $\mathbb{R}$ by hypothesis,
%                               and $\alpha$ functions as a lower bound for all
%                               $n\ge N$, $(p_{n})$ cannot have an upper bound. Thus, for any
%                               interval $(\alpha', \beta')$, where $\alpha<\alpha'<\beta'$,
%                               there must exist an $N'$ such that for all
%                               $n\ge N'$,  $p_{n}\ge \alpha'$. Since $\alpha'$ was arbitrary,
%                               this shows that $(p_{n})\to \infty$.
%                         \item There exists an $N$ such that for all $n\ge N$, $p_{n} \leq \beta$.

%                               Since $(p_{n})$ is unbounded in $\mathbb{R}$ by hypothesis,
%                               and $\beta$ functions as an upper bound for all
%                               $n\ge N$, $(p_{n})$ cannot have an lower bound. Thus, for any
%                               interval $(\alpha', \beta')$, where $\alpha'<\beta'<\beta$,
%                               there must exist an $N'$ such that for all $n\ge N'$,
%                               $p_{n}\le \beta'$. Since $\beta'$ was arbitrary, this shows that 
%                               $(p_{n})\to -\infty$.
%                   \end{enumerate}

%       \end{enumerate}


% \end{proof}

\end{document}